% GENERATED FILE. Edit canonical lessons, then run npm run generate:books.
% canonical-source-hash: fnv1a64:60c1f56430a8d320
% canonical-lessons: ML-C302-tunnuka, ML-C302-istiriyituka, ML-C302-murikkuka, ML-C302-ariyuka, ML-C302-araykkuka

\chapter{To sew, To iron, To cut, To chop, To grind}
\label{ch:ml-to-sew-to-iron-to-cut-to-chop-to-grind}
\hlchaptermodality{302}

\begin{chapteropening}
\textbf{By the end of this chapter:} \emph{I can say to sew, to iron, to cut, to chop and to grind.}

Complete the last lesson of chapter 302: I can say to sew, to iron, to cut, to chop and to grind.
\end{chapteropening}

\section[\texorpdfstring{\ml{തുന്നുക} (tunnuka)}{തുന്നുക (tunnuka)}]{\ml{തുന്നുക} (tunnuka) — to sew}

\label{lesson:ML-C302-tunnuka}



\begin{quote}
Before the new one: say the Malayalam for a stapler, then the Malayalam for a ruler.
\end{quote}

\subsection*{You'll want to know: \ml{തുന്നുക}}
\textbf{\ml{തുന്നുക}} — \emph{tunnuka} — \textquotedblleft{}to sew\textquotedblright{}.

\begin{grammarlens}[title={What you've built}]
One more word for this chapter.
\end{grammarlens}

\subsection*{Your turn}
\begin{itemize}
\raggedright
  \item \emph{Say it:} \emph{tunnuka}
  \item \emph{Say it:} \emph{tunnuka}, once more
  \item \emph{Say it:} say \emph{tunnuka}
  \item \emph{Recall:} say the Malayalam for a stapler, then the Malayalam for a ruler, then say \emph{tunnuka} again
\end{itemize}

\begin{tcolorbox}[breakable,colback=teal!4,colframe=teal!35!black,title={Before you move on}]
What does \ml{തുന്നുക} mean? (\textquotedblleft{}To sew\textquotedblright{}.) Say it once more.
\end{tcolorbox}
\section[\texorpdfstring{\ml{ഇസ്തിരിയിടുക} (istiriyiṭuka)}{ഇസ്തിരിയിടുക (istiriyiṭuka)}]{\ml{ഇസ്തിരിയിടുക} (istiriyiṭuka) — to iron (clothes)}

\label{lesson:ML-C302-istiriyituka}



\begin{quote}
Before the new one: say the Malayalam for a ruler, then the Malayalam for to sew.
\end{quote}

\subsection*{You'll want to know: \ml{ഇസ്തിരിയിടുക}}
\textbf{\ml{ഇസ്തിരിയിടുക}} — \emph{istiriyiṭuka} — \textquotedblleft{}to iron (clothes)\textquotedblright{}.

\begin{grammarlens}[title={What you've built}]
One more word for this chapter.
\end{grammarlens}

\subsection*{Your turn}
\begin{itemize}
\raggedright
  \item \emph{Say it:} \emph{istiriyiṭuka}
  \item \emph{Say it:} \emph{istiriyiṭuka}, once more
  \item \emph{Say it:} say \emph{istiriyiṭuka}
  \item \emph{Recall:} say the Malayalam for a ruler, then the Malayalam for to sew, then say \emph{istiriyiṭuka} again
\end{itemize}

\begin{tcolorbox}[breakable,colback=teal!4,colframe=teal!35!black,title={Before you move on}]
What does \ml{ഇസ്തിരിയിടുക} mean? (\textquotedblleft{}To iron (clothes)\textquotedblright{}.) Say it once more.
\end{tcolorbox}
\section[\texorpdfstring{\ml{മുറിക്കുക} (muṟikkuka)}{മുറിക്കുക (muṟikkuka)}]{\ml{മുറിക്കുക} (muṟikkuka) — to cut}

\label{lesson:ML-C302-murikkuka}



\begin{quote}
Before the new one: say the Malayalam for to sew, then the Malayalam for to iron.
\end{quote}

\subsection*{You'll want to know: \ml{മുറിക്കുക}}
\textbf{\ml{മുറിക്കുക}} — \emph{muṟikkuka} — \textquotedblleft{}to cut\textquotedblright{}.

\begin{grammarlens}[title={What you've built}]
One more word for this chapter.
\end{grammarlens}

\subsection*{Your turn}
\begin{itemize}
\raggedright
  \item \emph{Say it:} \emph{muṟikkuka}
  \item \emph{Say it:} \emph{muṟikkuka}, once more
  \item \emph{Say it:} say \emph{muṟikkuka}
  \item \emph{Recall:} say the Malayalam for to sew, then the Malayalam for to iron, then say \emph{muṟikkuka} again
\end{itemize}

\begin{tcolorbox}[breakable,colback=teal!4,colframe=teal!35!black,title={Before you move on}]
What does \ml{മുറിക്കുക} mean? (\textquotedblleft{}To cut\textquotedblright{}.) Say it once more.
\end{tcolorbox}
\section[\texorpdfstring{\ml{അരിയുക} (ariyuka)}{അരിയുക (ariyuka)}]{\ml{അരിയുക} (ariyuka) — to chop (vegetables)}

\label{lesson:ML-C302-ariyuka}



\begin{quote}
Before the new one: say the Malayalam for to iron, then the Malayalam for to cut.
\end{quote}

\subsection*{You'll want to know: \ml{അരിയുക}}
\textbf{\ml{അരിയുക}} — \emph{ariyuka} — \textquotedblleft{}to chop (vegetables)\textquotedblright{}.

\begin{grammarlens}[title={What you've built}]
One more word for this chapter.
\end{grammarlens}

\subsection*{Your turn}
\begin{itemize}
\raggedright
  \item \emph{Say it:} \emph{ariyuka}
  \item \emph{Say it:} \emph{ariyuka}, once more
  \item \emph{Say it:} say \emph{ariyuka}
  \item \emph{Recall:} say the Malayalam for to iron, then the Malayalam for to cut, then say \emph{ariyuka} again
\end{itemize}

\begin{tcolorbox}[breakable,colback=teal!4,colframe=teal!35!black,title={Before you move on}]
What does \ml{അരിയുക} mean? (\textquotedblleft{}To chop (vegetables)\textquotedblright{}.) Say it once more.
\end{tcolorbox}
\section[\texorpdfstring{\ml{അരയ്ക്കുക} (araykkuka)}{അരയ്ക്കുക (araykkuka)}]{\ml{അരയ്ക്കുക} (araykkuka) — to grind}

\label{lesson:ML-C302-araykkuka}



\begin{quote}
Before the new one: say the Malayalam for to cut, then the Malayalam for to chop.
\end{quote}

\subsection*{You'll want to know: \ml{അരയ്ക്കുക}}
\textbf{\ml{അരയ്ക്കുക}} — \emph{araykkuka} — \textquotedblleft{}to grind\textquotedblright{}.

\begin{grammarlens}[title={What you've built}]
One more word for this chapter.
\end{grammarlens}

\subsection*{Your turn}
\begin{itemize}
\raggedright
  \item \emph{Say it:} \emph{araykkuka}
  \item \emph{Say it:} \emph{araykkuka}, once more
  \item \emph{Say it:} say \emph{araykkuka}
  \item \emph{Recall:} say the Malayalam for to cut, then the Malayalam for to chop, then say \emph{araykkuka} again
\end{itemize}

\begin{tcolorbox}[breakable,colback=teal!4,colframe=teal!35!black,title={Before you move on}]
What does \ml{അരയ്ക്കുക} mean? (\textquotedblleft{}To grind\textquotedblright{}.) Say it once more.
\end{tcolorbox}
